\pg{335}b) \emph{O ente, que é contingente, consta de acto e} \textbf{\emph{de potencia.}} — O ente, que é \emph{contingente}, não existe pela exigencia da sua essencia. Ora o ente, que não existe pela exigencia da sua essencia, consta de existencia—\emph{realmente}—distincta—da essencia, e por isso de \emph{acto} e de \emph{potencia}. Logo o ente, que é \emph{contingente}, consta de \emph{acto} e de \emph{potencia} \footnote{(1) A essencia do ente \emph{contingente}, mesmo quando é actual e existe na realidade, distingue-se realmente da \emph{existencia}, como a potencia real se distingue do acto real. Portanto o ente \emph{contingente} consta de \emph{acto} e de \emph{potencia}. — D'ahi se segue que o ente \emph{contingente} não é necessariamente unico, mas podem existir muitos entes \emph{contingentes}. Por quanto, o ser do \emph{contingente}, sendo recebido na essencia, pode multiplicar-se indefinitamente, conforme a multidão dos sujeitos, ou das essencias, em que é recebido.}.

\textbf{58. O ente necessario é eterno.} — O ente \emph{necessario} é o que, pela exigencia da sua essencia, não pode deixar de existir. Ora o ente, que não pode deixar de existir, existiu sempre, existirá sempre, é \emph{eterno}. Logo o ente \emph{necessario} é \emph{eterno}.

\textbf{59. O ente contingente é produzido por uma causa extrinseca.} — Tudo o que se encontra num sujeito, e não pertence á sua essencia, deve ser produzido — ou pelos principios essenciaes do proprio sujeito — ou por uma causa extrinseca. Não ha meio termo. Ora a existencia, que não pertence á essencia do ente \emph{contingente}, não pode ser produzida pelos principios essenciaes do proprio ente; porque, neste caso, a essencia preexistiria á existencia (pois a causa preexiste ao effeito): o que é absurdo. Logo deve ser produzida por uma causa extrinseca. Logo o ente \emph{contingente} é produzido por uma causa extrinseca\footnote{(2) O ente \emph{necessario}, não podendo deixar de ser, é \emph{eterno}; porque, se não fosse eterno, se tivesse começo ou fim, seria \emph{contingente}. — Além d'isso, o ente \emph{necessario}, sendo \emph{acto purissimo}, não pode ter nem principio, nem fim, nem successão, porque tudo isto envolve potencialidade; e por isso é \emph{eterno}, porque a eternidade exclue principio, fim e successão. — Não só o ente \emph{necessario} é \emph{eterno}, mas é a sua \emph{eternidade}; porque elle é o seu ser, e a sua duração ou permanencia no \emph{ser} não se distingue da sua essencia, e por isso, assim como é a sua essencia, assim tambem é a sua eternidade. (\emph{Sum}. \emph{Th}. p. I, q. 10, a. 2; \emph{c. Gent.}, l. I, c. 15).\par (3) O ente \emph{contingente}, por não ser o proprio ser ou a propria existencia, isto é, por não ser \emph{necessario}, é \emph{indifferente} para a existencia e para a não-\pgn{336}existencia, e por isso não pode ser eterno, e teve começo. — E como recebeu elle a existencia real? Foi pela exigencia da sua natureza? Não. Porquanto, um sujeito, que é por si indifferente para dois attributos oppostos, não pode ter na sua essencia a razão pela qual exija um attributo, e não o opposto. Ora o ente \emph{contingente} é, pela sua essencia, indifferente para a existencia e para a não-existencia, é por isso a existencia não lhe pode convir pela mesma essencia. Logo o ente \emph{contingente} deve receber de um agente extrinseco a existencia real. (\emph{Sum}. \emph{Th}., p. I, q. 3, a. 4).}
